\subsection{Finite De Bruijn relations and the projective parity map}
\label{subsec:finite-debruijn-relations}

Laarhoven--de Weger's finite objects must first be typed as relations.  For
$k\geq1$ put
\[
 V_k=\Z/2^k\Z,
 \qquad
 \mathcal P(V_k)=\{A:A\subseteq V_k\}.
\]
Their modular Collatz graph has an edge $[a]_k\to[b]_k$ precisely when
integer lifts $a_1,b_1$ exist with $T(a_1)=b_1$.  Its forward relation is
therefore
\begin{equation}
\label{eq:LDW-finite-Collatz-relation}
 \mathcal T_k([a]_k)
 =\bigl\{[T(a)]_k,[T(a)+2^{k-1}]_k\bigr\}
 \quad:V_k\longrightarrow\mathcal P(V_k).
\end{equation}
Here $a$ may be any integer representative.  All representatives have the
same parity, and changing $a$ by $2^k$ changes its branch value by an odd
multiple of $2^{k-1}$.  Thus the two displayed residues are exactly all
possible outputs and are distinct.

Write a binary word least-significant digit first and identify
$b_0\cdots b_{k-1}$ with $\sum_{i<k}b_i2^i\in V_k$.  The binary De Bruijn
forward relation is
\begin{equation}
\label{eq:LDW-finite-shift-relation}
 \sigma_{2,k}(b_0\cdots b_{k-1})
 =\{b_1\cdots b_{k-1}c:c\in\{0,1\}\}.
\end{equation}
Numerically, this is the right shift together with that value plus
$2^{k-1}$.  Neither \eqref{eq:LDW-finite-Collatz-relation} nor
\eqref{eq:LDW-finite-shift-relation} is a deterministic self-map.

The source calls its finite encoder $\Phi_k$.  To prevent its dimension
subscript from colliding with Bernstein--Lagarias's multiplier subscript in
$\Phi_3=Q_3^{-1}$, this edition writes the same map as
\begin{equation}
\label{eq:LDW-finite-parity-map}
 \overline Q_{3,k}:V_k\longrightarrow V_k,
 \qquad
 \overline Q_{3,k}([a]_k)
 =\sum_{i=0}^{k-1}(T^i(a)\bmod2)2^i.
\end{equation}
Thus $\overline Q_{3,k}$ is source $\Phi_k$ without a normalization or a
change of orientation.  The symbol $3$ in $Q_3$ and $\Phi_3$ denotes the
$3x+1$ multiplier; it is not the graph dimension.

Everett's earlier Theorem~1 is exactly the bijectivity of this same finite
chronological map on the representatives $0\leq a<2^k$, together with the
stronger integer-fibre coordinates in
Theorem~\ref{thm:Everett-finite-fibres}.  Laarhoven--de Weger's later theorem
adds the two-successor graph relation; that relation is not retroactively
attributed to Everett.

For a point map $f:V_k\to V_k$ and $A\subseteq V_k$, write
$f[A]=\{f(a):a\in A\}$.  This elementwise extension is required in the next
statement.

\begin{theorem}[Laarhoven--de Weger finite graph theorem, reconstructed]
\label{thm:LDW-finite-debruijn}
For every $k\geq1$, $\overline Q_{3,k}$ is a bijection and
\begin{equation}
\label{eq:LDW-finite-relation-conjugacy}
 \overline Q_{3,k}[\mathcal T_k(a)]
 =\sigma_{2,k}(\overline Q_{3,k}(a))
 \qquad(a\in V_k).
\end{equation}
Consequently it is a directed-graph isomorphism from the binary modular
Collatz graph to $B(2,k)$.  Its inverse is the reduction modulo $2^k$ of
Bernstein--Lagarias's $\Phi_3=Q_3^{-1}$.
\end{theorem}

\begin{proof}
Proposition~\ref{prop:BL-Q} gives the exact congruence equivalence
\[
 x\equiv y\pmod{2^k}
 \quad\Longleftrightarrow\quad
 Q_3(x)\equiv Q_3(y)\pmod{2^k}.
\]
It proves that \eqref{eq:LDW-finite-parity-map} is well-defined and injective;
finiteness makes it bijective.  It also identifies it pointwise as the
reduction of $Q_3$, so reduction of $\Phi_3$ supplies the inverse.

Let $a_0$ and $a_0+2^k$ represent the two lift classes modulo $2^{k+1}$.
Their first $k$ parity digits agree, while their next parity digits differ,
again by the exact isometry.  Applying $T$ deletes the common first digit.
The two resulting length-$k$ itinerary words therefore have common first
$k-1$ digits and both possible terminal digits.  They are precisely
$\sigma_{2,k}(\overline Q_{3,k}([a_0]_k))$.  On the state side their residues
are exactly the two values in
\eqref{eq:LDW-finite-Collatz-relation}.  This proves
\eqref{eq:LDW-finite-relation-conjugacy} as equality of two-element subsets.
\end{proof}

The active source theorem has no proof environment.  Its preceding prose
depends on an earlier Lagarias congruence, and source line 131 names the next
digit as $x_{k+1}$ although the definitions force $x_k$.  A proof at source
lines 160--166 is commented out; its first sentence reverses the image edge,
although its following equivalence has the orientation used above.  These are
source defects, not alternative graph conventions
\cite[Sec.~2]{LaarhovenDeWeger2012}.

\begin{proposition}[exact path multiplicities]
\label{prop:LDW-path-counts}
Let $A_k$ be the adjacency matrix of either graph in
Theorem~\ref{thm:LDW-finite-debruijn}, and let $J$ be the $2^k$-by-$2^k$
all-ones matrix.  Then
\begin{equation}
\label{eq:LDW-path-counts}
 A_k^k=J,
 \qquad
 A_k^\ell=2^{\ell-k}J\quad(\ell\geq k).
\end{equation}
\end{proposition}

\begin{proof}
After $k$ De Bruijn steps all $k$ initial digits have been discarded.  For a
specified terminal word, exactly one ordered choice of the $k$ appended
digits reaches it.  After $\ell\geq k$ steps the first $\ell-k$ appended
digits are discarded and arbitrary, while the last $k$ are forced by the
terminal word.  This gives respectively $1$ and $2^{\ell-k}$ directed paths.
The graph isomorphism transports the count to the modular Collatz graph.
\end{proof}

\begin{nonclaim}
Equation \eqref{eq:LDW-path-counts} yields a uniform endpoint modulo $2^k$
only after a probability model is specified---for example, independent
uniform choices of the two outgoing finite-graph edges, equivalently uniform
unresolved higher input bits.  It does not make the deterministic integer
$T^k(n)$ random, and it does not prove that the known low bits are
statistically independent of an integer iterate.  The source's phrase that
they tell us ``absolutely nothing'' is not used without this qualification.
\end{nonclaim}

\subsection{Edge projections, the inverse limit, and literal integer edges}

Let $\rho_{k+1,k}:V_{k+1}\to V_k$ be reduction.  Directly from itinerary
prefixes,
\begin{equation}
\label{eq:LDW-projective-encoder}
 \rho_{k+1,k}\circ\overline Q_{3,k+1}
 =\overline Q_{3,k}\circ\rho_{k+1,k}.
\end{equation}
Let $E_k\subseteq V_k\times V_k$ be the edge relation of $\mathcal T_k$.
Division by two loses exactly one congruence digit, so
\begin{equation}
\label{eq:LDW-edge-lower-precision}
 ([a]_k,[b]_k)\in E_k
 \quad\Longleftrightarrow\quad
 b\equiv T(a)\pmod{2^{k-1}}.
\end{equation}
The right side is independent of the representative $a$ because its parity
and branch value modulo $2^{k-1}$ are fixed by $[a]_k$.

\begin{proposition}[exact graph inverse limit and natural-edge recovery]
\label{prop:LDW-edge-inverse-limit}
Reduction of both endpoints maps $E_{k+1}$ onto $E_k$.  The inverse limit of
the vertex systems is $\Ztwo$, and its compatible edge relation is exactly
the graph of the deterministic map $T:\Ztwo\to\Ztwo$.  Under
\eqref{eq:LDW-projective-encoder}, the inverse-limit graph isomorphism is
$Q_3$, and the target relation is the graph of the deterministic digit shift
$S$.

Moreover, for fixed literal $a,b\in\Nzero$, once both lie in
$\{0,\ldots,2^k-1\}$,
\begin{equation}
\label{eq:LDW-eventual-natural-edge}
 \begin{aligned}
 b=T(a)
 &\quad\Longleftrightarrow\quad
 ([a]_k,[b]_k)\in E_k\text{ for all sufficiently large }k,\\
 &\quad\Longleftrightarrow\quad
 ([a]_k,[b]_k)\in E_k\text{ for infinitely many }k.
 \end{aligned}
\end{equation}
\end{proposition}

\begin{proof}
Equation \eqref{eq:LDW-edge-lower-precision} shows that every higher edge
reduces to a lower edge.  Conversely, either lift of the initial vertex fixes
one more input bit and hence selects one of the two lower successors; choosing
that lift and either next higher bit supplies a higher edge.  Thus edge
projection is onto.  Notice that for one fixed higher vertex its two
successors reduce to a singleton, not to the entire lower successor set.

Suppose compatible $x,y\in\Ztwo$ form an edge at every level.  Then
\eqref{eq:LDW-edge-lower-precision} gives
$y\equiv T(x)\pmod{2^{k-1}}$ for every $k$, so $y=T(x)$.  The converse follows
by reduction of an actual edge.  The same argument for suffixes gives the
graph of $S$, while \eqref{eq:LDW-projective-encoder} and
Proposition~\ref{prop:BL-Q} identify the limiting map as $Q_3$.

If $b=T(a)$, every sufficiently large literal pair is an edge.  Conversely,
an occurrence in $E_k$ gives $b\equiv T(a)\pmod{2^{k-1}}$.  Infinitely many
occurrences include one with $2^{k-1}>|b-T(a)|$, when the congruence forces
equality in $\Z$.  This proves \eqref{eq:LDW-eventual-natural-edge}.
\end{proof}

This is the typed repair of the source phrase ``edges that appear in
infinitely many'' changing finite graphs.  It distinguishes the common
ambient pair $(a,b)\in\Nzero^2$ from compatible residue vertices and does not
call the set-valued $\mathcal T_k$ a deterministic quotient map.

\subsection{Periodic components, predecessor fibres, and source corrections}

\begin{proposition}[exact binary cycle count]
\label{prop:LDW-cycle-count}
Let $M_k$ be the number of components of $C(\Ztwo)$ containing a cycle of
least length $k$.  Then
\begin{equation}
\label{eq:LDW-cycle-count}
 2^k=\sum_{d\mid k}dM_d,
 \qquad
 M_k=\frac1k\sum_{d\mid k}\mu(d)2^{k/d}
 =\frac{2^k}{k}+O(2^{k/2}).
\end{equation}
In particular,
\[
 (M_k)_{k\geq1}=(2,1,2,3,6,9,\ldots).
\]
\end{proposition}

\begin{proof}
The $Q_3$ conjugacy identifies $T$-cycles with shift orbits of purely
periodic binary sequences.  The infinite periodic repetition of every
length-$k$ block has a unique least shift period $d\mid k$.  A shift orbit of
least period $d$ contains exactly $d$ distinct phases, so counting all
length-$k$ blocks gives the first identity.
M\"obius inversion gives the second.  In the error sum every proper divisor
contributes at most $2^{k/2}$, and the number of divisors is at most $k$, which
gives the displayed asymptotic.  Direct substitution gives the listed values.
\end{proof}

The source prints the correct divisor identity and M\"obius formula, but its
asymptotic changes the bound variable from $k$ to an unbound $n$.  Its printed
sequence $(1,2,1,2,3,6,\ldots)$ contradicts both the formula and its own two
fixed-point components; Proposition~\ref{prop:LDW-cycle-count} records the
forced correction.  Likewise, a cyclic De Bruijn sequence gives a Hamiltonian
cycle.  The source calls it a path while its parenthesized wrap-around digits
supply the closing edge \cite[Secs.~2 and 4]{LaarhovenDeWeger2012}.

\begin{proposition}[exact predecessor fibres]
\label{prop:LDW-predecessor-fibres}
For every $y\in\Ztwo$,
\begin{equation}
\label{eq:LDW-predecessor-fibres}
 T^{-1}(y)=\left\{2y,\frac{2y-1}{3}\right\},
 \qquad
 S^{-1}(y)=\{2y,2y+1\}.
\end{equation}
Both fibres have two distinct elements.  In particular, the forward Collatz
map is not a function automorphism on any weak component.
\end{proposition}

\begin{proof}
The point $2y$ is even and maps to $y$.  Since $3$ is a unit in $\Ztwo$,
$(2y-1)/3$ is a $2$-adic integer; it is odd and its odd branch also maps to
$y$.  The two branch equations give no other preimages.  Equality of the two
would imply $4y=-1$, impossible in $\Ztwo$.  The shift formula is immediate
from its deleted first digit.  Both predecessors of a point lie in its weak
component, proving the final assertion.
\end{proof}

Source line 241 says the forward mapping on an aperiodic component is an
automorphism.  Equation \eqref{eq:LDW-predecessor-fibres} contradicts that
sentence under the ordinary meaning of function automorphism.  The separate
source assertion that all aperiodic components are pairwise isomorphic is not
used downstream.

\subsection{Generalized maps: exact unit gate and scaling quotient}

For odd $a,b$, Laarhoven--de Weger define
\[
 T^{(a,b)}(n)=
 \begin{cases}
  n/2,&n\equiv0\pmod2,\\
  (an+b)/2,&n\equiv1\pmod2.
 \end{cases}
\]
At retained modulus $2$, this is precisely the Matthews presentation
\[
 (m_0,r_0)=(1,0),
 \qquad
 (m_1,r_1)=(a,-b).
\]
Oddness of $a,b$ supplies the branch congruence and makes both multipliers
units.  Hence the finite and infinite parity encoders in source Section~5.1
are the $d=2$ specialization of
Theorem~\ref{thm:general-d-itinerary-conjugacy}; their orientation is
$\operatorname{It}_{T^{(a,b)}}$, not its inverse.  Only at
$(a,b)=(3,1)$ does this encoder specialize to $Q_3$.

For a retained base $d\geq2$, consider instead
\begin{equation}
\label{eq:LDW-general-p-branch}
 F(n)=\frac{a_i n+b_i}{d}
 \qquad(n\equiv i\pmod d).
\end{equation}
The exact hypotheses for the full digit-shift conclusion are
\begin{equation}
\label{eq:LDW-general-p-gate}
 a_i i+b_i\equiv0\pmod d,
 \qquad
 \gcd(a_i,d)=1
 \quad(0\leq i<d).
\end{equation}
The first is branch integrality.  For $n=i+dt$, the next digit is
\begin{equation}
\label{eq:LDW-next-digit-affine}
 F(i+dt)\equiv \frac{a_i i+b_i}{d}+a_it\pmod d.
\end{equation}
The second condition makes this an affine permutation of all $d$ next digits
with singleton fibres.  With $m_i=a_i$ and $r_i=-b_i$, conditions
\eqref{eq:LDW-general-p-gate} are exactly the retained Matthews congruence and
unit gate.  Theorem~\ref{thm:general-d-itinerary-conjugacy} therefore proves
the full finite and $d$-adic De Bruijn conclusions.  If a multiplier is not a
unit, Proposition~\ref{prop:nonunit-itinerary-obstruction} gives the missing
digits and non-singleton fibres explicitly.

The source's Section~5.2 says only ``appropriately chosen'' coefficients and
does not print \eqref{eq:LDW-general-p-gate}.  If that phrase means merely
integer-valued branches, choose $a_i=d,b_i=0$ for every $i$.  Then
$F=\operatorname{id}_{\Z}$, whose finite graph is not a De Bruijn graph.  The
source also leaves $x_i^{(f)}$ undefined and does not say that $p$ is prime
despite writing $\Z_p$.  The edition admits the claim only at the exact scope
above; for prime $d=p$ it is the usual $p$-adic statement.  The source's
ternary example has multipliers $2,4,4$, all units modulo $3$, and lies inside
the repaired scope.

\begin{proposition}[rational cycles and primitive $3n+b$ data]
\label{prop:LDW-rational-cycle-scaling}
Let $(x_0,\ldots,x_{L-1})$ be a rational $T$-cycle in $\Ztwo$, with cyclic
indices, and let $b>0$ be a common odd denominator.  Put $n_j=bx_j$.  Then
$(n_j)$ is an integer cycle of $T^{(3,b)}$, and
\begin{equation}
\label{eq:LDW-rational-scaling}
 T(x_j)=\frac1bT^{(3,b)}(n_j).
\end{equation}
Conversely every such integer cycle scales to a rational $T$-cycle.  The
resulting bijection is between rational cycles and the equivalence classes
generated by the scalings
\begin{equation}
\label{eq:LDW-scaling-equivalence}
 (b,(n_j))\sim(cb,(cn_j))
 \qquad(c\in\Npos\text{ odd}),
\end{equation}
or, equivalently, their unique representatives with $b>0$ and
$\gcd(b,n_0,\ldots,n_{L-1})=1$.  For a rational $T$-cycle this primitive
$b$ is automatically coprime to $3$.
\end{proposition}

\begin{proof}
Multiplication by odd $b$ preserves parity in $\Ztwo$.  On an even phase,
$bT(x_j)=n_j/2$.  On an odd phase,
$bT(x_j)=(3n_j+b)/2$.  This proves
\eqref{eq:LDW-rational-scaling} and both cycle directions.  Scaling numerator
and denominator by a common odd positive integer leaves every $x_j$ fixed.
Dividing by the common gcd terminates at a primitive representative, and any
two primitive positive-denominator representatives of the same rational
tuple are equal.  If the cycle has return length $L$ and contains $m$ odd
steps, the fixed-return equation shows that the reduced denominator of every
phase divides $2^L-3^m$.  Hence their primitive common denominator also
divides $2^L-3^m$.  If $m\geq1$, this integer is coprime to $3$.  If $m=0$,
every phase uses $x\mapsto x/2$, so periodicity forces the zero cycle and the
primitive common denominator is $1$.  Thus it is coprime to $3$ in both
cases.
\end{proof}

The source's first, unqualified ``one-to-one'' sentence requires
\eqref{eq:LDW-scaling-equivalence}: for every odd $b$, the integer cycle
$\{b,2b\}$ of $T^{(3,b)}$ scales to the same rational cycle $\{1,2\}$.
Its subsequent phrase ``having denominator $b$'' can instead be read as the
exact reduced common denominator; on that reading it already selects the
primitive representative and is correct.  Lagarias's primary treatment
explicitly retains the raw gcd strata and then defines this primitive layer
\cite[pp.~39--40]{Lagarias1990}.  The proposition makes both scopes explicit
without weakening the scaling identity or suppressing the raw scaling
fibres.

\begin{sourceissue}[remaining prose and hypothesis boundaries]
\label{issue:LDW-source-boundaries}
The source introduction describes absence of divergent paths by
$\lim_kT^k(n)<\infty$, although the $1\leftrightarrow2$ orbit has no limit;
the boundedness condition is $\sup_kT^k(n)<\infty$.  Source line 79 says
``for all $n$'' where its fixed-input argument requires ``for all $k$''.
The rational-periodicity statement in Sections~4 is explicitly a conjecture.
The average-multiplier discussion and the predicted $5n+1$ components are
heuristics; no theorem is attributed to either discussion.  Finally, the source's positive-integer
reformulation through $\frac13\Z\setminus\Z$ depends on earlier cited
literature and is not enlarged here before that exact dependency is checked.
\end{sourceissue}

\begin{nonclaim}
The finite graph theorem, exact path counts, inverse limit, cycle enumeration,
and generalized unit-gated conjugacy do not characterize
$Q_3(\Npos)\subset\Ztwo$.  They therefore do not prove that a positive
integer reaches $1$.  The arithmetic difficulty lies in the embedded subset
and its labels, not in the existence of the ambient De Bruijn isomorphism.
\end{nonclaim}
